\documentclass[11pt]{article}
\usepackage[a4paper,margin=27mm]{geometry}
\usepackage{amsmath,amssymb,amsthm,mathtools}
\usepackage{booktabs,longtable,array,enumitem,microtype}
\usepackage{xcolor}
\usepackage[hidelinks]{hyperref}
\usepackage[T1]{fontenc}
\usepackage{lmodern}

\definecolor{proved}{RGB}{0,105,75}
\definecolor{conditional}{RGB}{160,92,0}
\definecolor{open}{RGB}{165,25,35}
\newcommand{\Pstatus}{\textcolor{proved}{\textsf{PROVED}}}
\newcommand{\Cstatus}{\textcolor{conditional}{\textsf{CONDITIONAL}}}
\newcommand{\Ostatus}{\textcolor{open}{\textsf{OPEN}}}
\newcommand{\Rstatus}{\textcolor{proved}{\textsf{RESTRICTED}}}
\newcommand{\Aphys}{\mathcal A_{\mathrm{phys}}}
\newcommand{\R}{\mathsf R}
\newcommand{\B}{\mathfrak B_{\mathrm{rig}}}
\newcommand{\spec}{\operatorname{spec}}
\newcommand{\Cov}{\operatorname{Cov}}
\newcommand{\supp}{\operatorname{supp}}
\newcommand{\dist}{\operatorname{dist}}
\newcommand{\e}{\mathrm e}

\newtheorem{definition}{Definition}[section]
\newtheorem{theorem}[definition]{Theorem}
\newtheorem{proposition}[definition]{Proposition}
\newtheorem{lemma}[definition]{Lemma}
\newtheorem{corollary}[definition]{Corollary}
\newtheorem{counterexample}[definition]{Counterexample}
\theoremstyle{remark}
\newtheorem{remark}[definition]{Remark}
\newtheorem{obligation}[definition]{Frontier obligation}

\title{Representational Dynamics and Infrared Rigidity:\\
A Constructive Research Programme for the Yang--Mills Mass Gap}
\author{Research dossier}
\date{20 July 2026}

\begin{document}
\maketitle

\begin{abstract}
We formulate a representation-aware programme for the four-dimensional
Yang--Mills existence and mass-gap problem.  The physical theory is treated as
coherent information transported among regulated measures, observable
algebras, renormalization-group descriptions, Euclidean correlation
functionals, and reconstructed spectral data.  We define uniform physical
relational rigidity, prove a positive spectral transfer theorem, give an
abstract noncircular influence-basin theorem, type the required representation
morphisms, and exhibit counterexamples that invalidate several common
shortcuts.  We then isolate the first genuinely open four-dimensional step as
a cutoff- and volume-uniform full-measure basin-entry inequality.  The result
is a rigorous conditional architecture and a falsifiable research programme,
not a solution of the Clay problem.  Its strongest claim level is C3: a
conditional continuum/transfer theorem.
\end{abstract}

\tableofcontents

\section{Epistemic statement and exact target}

The Clay problem asks for a nontrivial quantum Yang--Mills theory on
$\mathbb R^4$ for every compact simple gauge group $G$, satisfying axioms at
least as strong as those conventionally required for a relativistic quantum
field theory, and possessing a mass gap $\Delta>0$ above the vacuum
\cite{JaffeWitten}.  This dossier does not prove that theorem.  It separates
the theorem into interfaces so that a proof cannot move spectral information
through a representation change without declaring what is preserved.

\begin{definition}[Clay target]
For a compact simple Lie group $G$, a \emph{Clay-level solution} supplies a
physical Hilbert space $\mathcal H_G$, a unit vacuum $\Omega_G$, a positive
self-adjoint Hamiltonian $H_G$, a nontrivial gauge-invariant local quantum
field theory satisfying the required relativistic axioms, and a number
$\Delta_G>0$ such that
\[
  \spec(H_G)\cap(0,\Delta_G)=\varnothing,
\]
with nonvacuum finite-energy content.  The construction must be obtained from
mathematically controlled data, not by postulating the Hamiltonian or gap.
\end{definition}

The mass gap is distinct from confinement, an area law, a one-particle pole,
or a gap for an auxiliary stochastic generator.  Any one of these may be
related to a physical gap under additional hypotheses; none is a synonym.

\subsection{Claim hierarchy}

\begin{center}
\begin{tabular}{cl}
\toprule
Level & Meaning\\
\midrule
C0 & concepts and definitions\\
C1 & finite or executable calibration\\
C2 & rigorous restricted-model theorem\\
C3 & conditional continuum or representation-transfer theorem\\
C4 & regulator-specific nonperturbative construction, uniform limits incomplete\\
C5 & continuum construction and gap for a restricted target\\
C6 & the full Clay target\\
\bottomrule
\end{tabular}
\end{center}

Nothing proved below exceeds C3.  This ceiling is a mathematical conclusion of
the dependency audit, not a rhetorical disclaimer.

\section{Ontology: a universe as coherent representational dynamics}

The proposed ontology does not make representation more real than physics.
It says instead that physically meaningful identity is the invariant content
of a coherently related family of partial descriptions.  A universe is
represented by
\[
 \mathfrak U=(\mathbf{Rep},\mathbf{Mor},\mathcal I,\mathcal P,\mathcal L),
\]
where $\mathbf{Rep}$ is a class of descriptions, $\mathbf{Mor}$ contains
admissible transformations among descriptions, $\mathcal I$ assigns invariant
physical content, $\mathcal P$ records preservation certificates, and
$\mathcal L$ specifies limiting coherence.  No individual regulator is the
universe.  The universe is the compatible structure that remains identifiable
through admissible changes of resolution and language.

In this interpretation:
\begin{itemize}[itemsep=2pt]
\item gauge potentials are coordinates on relational physical information;
\item regulators are temporary realizations that expose different proof
      structures;
\item renormalization is a controlled transformation of descriptive scale;
\item particles are spectral interpretations in a reconstructed
      representation;
\item the continuum theory is a coherence condition across limiting
      representations; and
\item a mass gap is invariant infrared rigidity: sufficiently low positive
      energy cannot be created by changing a faithful physical description.
\end{itemize}

The ontology becomes mathematical only after the words ``controlled'',
``faithful'', and ``invariant'' are replaced by typed maps, topologies, error
bounds, and density theorems.

\section{Physical observables and completeness}

Let $\Lambda_{a,L}$ be a regulated spacetime with cutoff $a$ and physical
linear size $L$.  Its kinematical fields form $\mathcal K_{a,L}$ and its gauge
group is $\mathcal G_{a,L}$.  The primary physical algebra is
\[
 \mathcal A_{\mathrm{phys},a,L}
  =\{F:\mathcal K_{a,L}\to\mathbb C:
          F(g\cdot A)=F(A)\text{ for all }g\in\mathcal G_{a,L}\},
\]
completed only after a state and topology have been specified.

Candidate cores include local curvature composites after renormalization,
smeared gauge-invariant operators, Wilson loops, and extended operators needed
to distinguish global or topological sectors.  Wilson loops alone are not
assumed complete.

\begin{definition}[State-separating core]
A $*$-subalgebra $\mathcal D\subset\Aphys$ is state separating if equality of
all expectation values on $\mathcal D$ implies equality of physical states.
After OS reconstruction it is \emph{spectrally complete} if
\[
 \overline{\{A\Omega:A\in\mathcal D\}}=\mathcal H.
\]
For a mass-gap argument it suffices to require density of centered vectors in
$\Omega^\perp$.
\end{definition}

\begin{obligation}[Observable completeness]
Prove that the regulated observable cores converge to a continuum core whose
centered vectors are dense in the reconstructed vacuum complement.  Until
then, any correlation result is sectorial.
\end{obligation}

\section{Representation packages and typed morphisms}

\begin{definition}[Representation package]
A representation is a tuple
\[
 \R=(\mathcal K,\Aphys,\mathcal S,\mathcal D,\mathcal G,
       \mathcal L,\mathcal Q),
\]
consisting of kinematics, physical observables, states or measures, dynamics,
gauge structure, locality/scale data, and quantitative control data.
\end{definition}

A morphism $F:\R\to\R'$ contains maps on states and observables with the dual
compatibility relation
\[
 \langle F_{\!*}\omega,A'\rangle
   \simeq \langle\omega,F^*A'\rangle,
\]
where $\simeq$ is equality for exact maps and an explicitly bounded error for
approximate maps.  Every morphism carries a certificate
\[
 \operatorname{cert}(F)=(g,p,\ell,o,u),
\]
recording gauge compatibility, reflection positivity, locality, observable
completeness, and uniformity in volume.

\begin{proposition}[Composition ledger \Pstatus]
Let $F:\R_0\to\R_1$ and $G:\R_1\to\R_2$ have additive expectation errors
$\varepsilon_F$ and $\varepsilon_G$.  Then $G\circ F$ has error at most
$\varepsilon_F+\varepsilon_G$, and
\[
 \operatorname{cert}(G\circ F)
 =\operatorname{cert}(F)\wedge\operatorname{cert}(G)
\]
componentwise.  Thus a lost certificate cannot reappear without a separate
recovery theorem.
\end{proposition}

\begin{proof}
Insert and subtract the intermediate expectation and use the triangle
inequality.  A property claimed for the composition must hold at both stages;
otherwise the corresponding Boolean certificate is false.  The executable
artifact checks this rule with exact rational errors.
\end{proof}

The important morphism classes are exact equivalence, regulator embedding,
coarse graining, renormalized observable transport, gauge fixing with global
residual data, reconstruction, and continuum limiting.  Calling every one of
them an ``equivalence'' would erase precisely the information the programme is
designed to audit.

\section{Regulator atlas}

The programme is not founded on earlier Wilson-action work.  The Wilson action
appears as one regulator among several because any final continuum conclusion
needs a regulator choice or a regulator-comparison theorem.

\begin{longtable}{p{.17\textwidth}p{.18\textwidth}p{.18\textwidth}p{.33\textwidth}}
\toprule
Regulator & Gauge invariance & Positivity/locality & Principal proof obligations\\
\midrule
\endhead
Wilson plaquette & exact lattice gauge invariance & local; standard reflection-positive constructions exist in suitable settings & weak-coupling continuum trajectory, full-measure RG, physical scaling, universality\\
Heat-kernel plaquette & exact class-function gauge invariance & local; positivity must be checked for the chosen time slicing & coefficient control, comparison with continuum renormalization, universality\\
Villain-type or auxiliary-field formulations & gauge invariance depends on nonabelian implementation & may expose convexity but positivity is formulation dependent & auxiliary-field erasure, topology, equivalence to target theory\\
Hamiltonian lattice gauge theory & Gauss law and physical Hilbert space explicit & spectral interpretation native; Euclidean covariance indirect & continuum limit, locality, Lorentz covariance, comparison with Euclidean measures\\
Gradient-flowed observables over another regulator & gauge-covariant smoothing of observables & improves ultraviolet observables but is not a stand-alone theory & invertibility/sufficiency of flowed data, zero-flow-time limit\\
\bottomrule
\end{longtable}

Balaban's work establishes deep multiscale control for lattice gauge theories,
including a four-dimensional small-field coupling renormalization analysis
\cite{BalabanI}.  It is indispensable evidence for the feasibility of exact
RG methods, but it does not supply the complete full-measure, continuum,
OS-reconstruction, and spectral chain required here.

\section{Uniform physical relational rigidity}

Let $a_n\downarrow0$, $L_n\uparrow\infty$, and let $\omega_n$ be regulated
states.  Let $\tau_t^{(n)}$ denote Euclidean time translation in physical
units, with interpolation when $t$ is not a lattice multiple.  For a centered
observable $A_n$, write
\[
 C_n(A_n;t)=\omega_n(A_n^*\tau_t^{(n)}A_n)
            -|\omega_n(A_n)|^2.
\]

\begin{definition}[Uniform physical relational rigidity]
The family has UPRR with rate $m>0$ on a continuum core $\mathcal D$ if every
$A\in\mathcal D$ has regulated representatives $A_n$, a constant $K_A<\infty$
independent of $n$ and $L_n$, and errors $r_n(A,t)$ such that
\begin{align}
 0\le C_n(A_n;t)&\le K_A\e^{-mt}+r_n(A,t),\label{eq:uprr}\\
 \lim_{n\to\infty}r_n(A,t)&=0\quad\text{for every fixed }t\ge0,\nonumber
\end{align}
and the quadratic forms converge to continuum correlations.  The parameter
$m$ has units of inverse physical length.  Polarization supplies mixed
correlators when required.
\end{definition}

The definition is noncircular only if \eqref{eq:uprr} is obtained from an
independent local or multiscale criterion.  It must not be used as the
definition of the RG basin whose entry one then claims to prove.

For a one-step transfer eigenvalue $\lambda_a$, the dimensionless gap is
$\delta_a=-\log\lambda_a$ while the physical energy is $\delta_a/a$.  A lower
bound on $\delta_a$ that shrinks faster than $a$ is irrelevant to the continuum
mass gap.

\section{Positive spectral transfer}

\begin{lemma}[Positive Laplace support lemma \Pstatus]\label{lem:laplace}
Let $\mu$ be a finite positive measure on $[0,\infty)$ and
\[
 C(t)=\int_{[0,\infty)}\e^{-t\lambda}\,d\mu(\lambda).
\]
If $C(t)\le K\e^{-mt}$ for all $t\ge0$, then
$\mu([0,m))=0$.
\end{lemma}

\begin{proof}
For $0<\varepsilon<m$,
\[
 \mu([0,m-\varepsilon])\e^{-t(m-\varepsilon)}
 \le C(t)\le K\e^{-mt}.
\]
Thus $\mu([0,m-\varepsilon])\le K\e^{-\varepsilon t}$ for every
$t$, and the left side is zero.  Let $\varepsilon\downarrow0$ through a
countable sequence.
\end{proof}

\begin{theorem}[Dense-core rigidity transfer \Pstatus]\label{thm:transfer}
Let $H\ge0$ be self-adjoint on $\mathcal H$, with unique vacuum $\Omega$, and
let $\mathcal D_0\subset\Omega^\perp$ be dense.  Suppose for every
$\psi\in\mathcal D_0$ there is $K_\psi<\infty$ such that
\[
 \langle\psi,\e^{-tH}\psi\rangle
 \le K_\psi\e^{-mt}\qquad(t\ge0).
\]
Then
\[
 \spec(H|_{\Omega^\perp})\subset[m,\infty).
\]
\end{theorem}

\begin{proof}
For $\psi\in\mathcal D_0$, the spectral theorem gives a positive measure
$\mu_\psi$ and Lemma~\ref{lem:laplace} shows
$\|E_H([0,m))\psi\|^2=0$.  The spectral projection is bounded, so density gives
$E_H([0,m))=0$ on $\Omega^\perp$.
\end{proof}

\begin{theorem}[Regulator-to-continuum transfer \Cstatus]\label{thm:regcont}
Assume regulated representatives of a continuum observable core satisfy UPRR,
their centered quadratic forms converge for every fixed $t$, reflection
positivity reconstructs $\e^{-tH}$, and their continuum vectors are dense in
$\Omega^\perp$.  Then $H$ has no spectrum in $(0,m)$.
\end{theorem}

\begin{proof}
Take $n\to\infty$ in \eqref{eq:uprr} at fixed $t$ to obtain the hypothesis of
Theorem~\ref{thm:transfer} on the dense continuum core.  Apply that theorem.
\end{proof}

Fixed-time error convergence is enough for this conditional theorem because
the continuum estimate is established separately for every $t$.  It is not
enough to infer a uniform regulated spectral gap if the large-time and cutoff
limits are interchanged before a continuum semigroup exists.

\section{A noncircular infrared basin}

We now give an abstract basin whose membership is local in a block
specification rather than defined by the desired global correlation decay.
Let $I$ index blocks with metric $d$.  For a local observable $F$, let
$\delta_i(F)$ be its oscillation under a change of block $i$.  Let
$C=(c_{ij})$ be a nonnegative influence matrix derived from conditional kernels.

\begin{definition}[Weighted influence basin]
For $\kappa>0$ and $q<1$, an effective theory lies in
$\B(\kappa,q)$ if its conditional specification obeys the comparison estimate
\begin{equation}\label{eq:comparison}
 |\Cov(F,G)|
 \le K_0\sum_{i,j}\delta_i(F)D_{ij}\delta_j(G),
 \qquad D=(I-C)^{-1},
\end{equation}
and
\begin{equation}\label{eq:weighted}
 \|C\|_\kappa
 :=\sup_i\sum_j \e^{\kappa d(i,j)}c_{ij}\le q<1.
\end{equation}
All constants must be uniform in the total volume.  In a concrete Gibbs
specification, \eqref{eq:comparison} is to be proved by a Dobrushin comparison
argument; it is displayed rather than hidden inside the word ``mixing''.
\end{definition}

\begin{theorem}[Basin-to-clustering \Cstatus]\label{thm:basin}
If an effective theory lies in $\B(\kappa,q)$ and $F,G$ have finite block
supports $S_F,S_G$, then
\[
 |\Cov(F,G)|
 \le \frac{K_0}{1-q}\,
      \e^{-\kappa d(S_F,S_G)}
      \Big(\sum_i\delta_i(F)\Big)
      \Big(\sum_j\delta_j(G)\Big).
\]
\end{theorem}

\begin{proof}
The triangle inequality implies that the weighted row norm is submultiplicative,
so $D=\sum_{n\ge0}C^n$ and $\|D\|_\kappa\le(1-q)^{-1}$.
Consequently $D_{ij}\le(1-q)^{-1}\e^{-\kappa d(i,j)}$.
Insert this in \eqref{eq:comparison} and use
$d(i,j)\ge d(S_F,S_G)$ for nonzero oscillations.
\end{proof}

\begin{corollary}[RG scale to physical rigidity \Cstatus]
Suppose a gauge-compatible RG flow reaches $\B(\kappa,q)$ at block length
$\ell_*$; observable pullback has summable errors and uniformly controlled
support growth.  Then the original correlations decay with physical rate at
least
\[
 m_{\rm corr}=\frac{\kappa}{\ell_*},
\]
up to the declared pullback errors.  If the controlled direction is Euclidean
time and OS hypotheses hold, Theorem~\ref{thm:regcont} yields the same lower
bound for the Hamiltonian gap.
\end{corollary}

This result separates two questions: the implication from basin membership is
proved, while membership of the four-dimensional Yang--Mills trajectory is not.

\section{Gauge-covariant RG and the basin-entry frontier}

A complete RG step must transform the measure, action, observables, scale, and
reconstruction data together.  Schematically,
\[
 (\mu_k,S_k,\mathcal A_{\mathrm{phys},k},Z_k,\ell_k)
 \longmapsto
 (\mu_{k+1},S_{k+1},\mathcal A_{\mathrm{phys},k+1},Z_{k+1},b\ell_k).
\]
It must report gauge covariance, locality or quasilocality, reflection
positivity, volume uniformity, running couplings, irrelevant remainders, and
errors.  A perturbative beta function alone is not this map.

The large-field decomposition must be gauge invariant, for example through
plaquette-curvature or orbit-distance events.  Exceptional blocks require
activities summable over positions, polymers, and RG scales.  Global holonomy,
center flux, nontrivial bundles, and instanton-like sectors must either be
included or bounded.  A coordinate chart near the trivial connection does not
cover these obligations.

\begin{definition}[Basin-entry statement $\mathrm{BE}$]
$\mathrm{BE}(k_*,\kappa,q)$ is the assertion that, along the cutoff trajectory
defining four-dimensional pure Yang--Mills, there is a finite physical entry
scale $\ell_*=b^{k_*}a$ such that the complete effective conditional kernels
$\gamma_B^{(k_*)}$ have influence coefficients satisfying
\begin{equation}\label{eq:BE}
 \sup_B\sum_C \e^{\kappa d(B,C)}c_{BC}^{(k_*)}\le q<1,
\end{equation}
and satisfy the comparison estimate \eqref{eq:comparison}, with all constants
uniform in volume and cutoff and all reconstruction errors summable.
\end{definition}

\begin{obligation}[Central frontier \Ostatus]\label{obl:BE}
Prove $\mathrm{BE}(k_*,\kappa,q)$ for the full four-dimensional pure
Yang--Mills measure, or prove an equally explicit noncircular polymer,
coercivity, or transfer-contraction criterion.  No such proof is supplied here.
\end{obligation}

Asymptotic freedom controls the ultraviolet running of the coupling but does
not imply \eqref{eq:BE}.  Strong-coupling lattice theorems, such as exponential
correlation decay established by functional-inequality methods in explicit
parameter regimes \cite{SZZ,Nissim}, are rigorous restricted calibrations.
They do not by themselves connect the weak-coupling continuum trajectory to a
positive physical gap.

\subsection{Concrete research sub-obligations}

\begin{enumerate}[label=\textbf{BE\arabic*.},leftmargin=3em]
\item Derive a gauge-orbit block specification with quantitatively local
      conditional kernels.
\item Establish weak-coupling small-field bounds with constants stable for the
      required number of RG steps.
\item Prove a summable exceptional-polymer estimate uniform in volume.
\item Cover global orbit and topological sectors and control atlas overlaps.
\item Bound generated nonlocal interactions in a weighted norm strong enough
      for \eqref{eq:BE}.
\item Show that the flow enters $q<1$ at finite physical $\ell_*$ rather than
      merely at a cutoff-dependent scale.
\item Preserve or recover reflection positivity and retain a spectrally
      complete observable transport.
\end{enumerate}

These are detailed todos in the literal mathematical sense: each names a
finite theorem or inequality whose proof can be accepted or rejected.

\section{Continuum coherence}

Let $S_n^{(a,L)}$ denote renormalized $n$-point Schwinger functionals on a
gauge-invariant test algebra.  A continuum theory requires more than a formal
limit of actions.

\begin{definition}[Continuum coherence package]
A package consists of:
\begin{enumerate}[label=(\roman*)]
\item a topology in which all $S_n^{(a,L)}$ converge;
\item Euclidean covariance, symmetry, locality, reflection positivity, and
      the regularity conditions required by the chosen OS theorem;
\item nontriviality and a certificate identifying the limit as pure
      Yang--Mills under stated renormalization conditions;
\item uniqueness or a regulator-comparison theorem in a declared class; and
\item justified interchange of cutoff, volume, RG-depth, observable, and
      large-separation limits.
\end{enumerate}
\end{definition}

The original OS programme reconstructs relativistic quantum theory from
suitable Euclidean Green functions \cite{OS1,OS2}.  The precise regularity
hypotheses matter: the corrected second paper is not interchangeable with an
informal list of covariance and positivity properties.

\begin{theorem}[Conditional continuum gap theorem \Cstatus]
Assume a continuum coherence package for nontrivial pure Yang--Mills, UPRR with
rate $m>0$ on an OS-dense gauge-invariant core, a unique vacuum, and nonempty
finite-energy nonvacuum content.  Then the reconstructed Hamiltonian obeys
\[
 \spec(H)=\{0\}\cup\Sigma,
 \qquad \Sigma\subset[m,\infty),\quad\Sigma\ne\varnothing.
\]
\end{theorem}

\begin{proof}
The coherence package and OS theorem construct $\mathcal H,\Omega,H$.  UPRR
passes to the continuum by Theorem~\ref{thm:regcont}.  Density upgrades the
sectorial conclusion to all of $\Omega^\perp$, and the final assumption gives
$\Sigma\ne\varnothing$.
\end{proof}

\begin{obligation}[Continuum construction \Ostatus]
Construct the above package in four dimensions with uniform estimates strong
enough to pass UPRR, and prove regulator independence.  This dossier does not
solve that constructive field-theory problem.
\end{obligation}

\section{Counterexample bank}

\begin{counterexample}[Finite-volume fake gap]
Let $H_L=\operatorname{diag}(0,L^{-1})$.  Every finite $L$ has a positive gap,
but the gaps converge to zero.  Positivity without volume uniformity proves no
infinite-volume gap.
\end{counterexample}

\begin{counterexample}[Lattice-unit collapse]
Let $a_n=n^{-1}$ and let a transfer operator have dimensionless gap
$\delta_n=a_n^2$.  Then $\delta_n>0$ for every regulator, while the physical
energy is $\delta_n/a_n=a_n\to0$.
\end{counterexample}

\begin{counterexample}[Hidden gapless sector]
Take
\[
 \mathcal H=\mathbb C\Omega\oplus\mathbb C g\oplus L^2([0,1]),
 \qquad H=0\oplus1\oplus M_\lambda.
\]
An observable class whose centered vectors span only $\mathbb Cg$ has exact
decay $\e^{-t}$, yet $H$ has spectrum arbitrarily close to zero in the hidden
$L^2$ sector.  Observable density is indispensable.
\end{counterexample}

\begin{counterexample}[Why positivity matters]
If signed spectral weights are permitted, low-energy terms can cancel.  For
example $\delta_{\lambda}-\delta_{\lambda}$ has zero Laplace transform while
being formally supported at any chosen low energy.  Reflection positivity is
what turns Euclidean quadratic correlations into positive spectral measures.
\end{counterexample}

\begin{counterexample}[Stochastic gap is not physical energy]
For a Markov generator $L$ with invariant measure $\mu$, replacing $L$ by
$cL$ leaves $\mu$ unchanged and multiplies its Poincar\'e spectral gap by $c$.
The Euclidean quantum theory determined by $\mu$ has not thereby acquired an
arbitrarily large physical mass.  A comparison theorem must link the sampling
dynamics to physical time before its mixing rate has spectral meaning.
\end{counterexample}

\begin{counterexample}[A vacuous local sector]
A topological theory or an observable subalgebra with only multiples of the
identity can satisfy any centered local correlation estimate vacuously.  The
Clay target also requires a nontrivial theory and physical nonvacuum content.
\end{counterexample}

\section{Calibration matrix}

\begin{center}
\small
\begin{tabular}{p{.22\textwidth}p{.16\textwidth}p{.18\textwidth}p{.15\textwidth}p{.16\textwidth}}
\toprule
Model & UPRR & Complete core & Physical gap & Verdict\\
\midrule
finite diagonal $H=\mathrm{diag}(0,1,2)$ & yes & yes & yes & accepts\\
massive Gaussian field & yes & standard field core & yes & accepts\\
massless Gaussian field & no uniform $m>0$ & yes & no & rejects\\
$H_L=\mathrm{diag}(0,L^{-1})$ & nonuniform & yes & no limit gap & rejects\\
hidden gapless direct sum & sectorial yes & no & no & rejects\\
topological/trivial local sector & vacuous & generally no & not sufficient & rejects\\
strong-coupling lattice YM & restricted decay & regulator core & lattice result & calibration only\\
four-dimensional continuum YM & desired & open & open & frontier\\
\bottomrule
\end{tabular}
\end{center}

The finite diagonal, nonuniform, hidden-sector, positivity, and morphism-error
cases are implemented in \texttt{finite\_spectral\_core.py} and its unit tests.
Because no proof assistant is installed in the execution environment, the
artifact uses exact rational arithmetic and fail-closed runtime validation.
It is an executable C1 certificate, not formal verification of the analytic
frontier.

\section{Limit-order ledger}

The following implications require named interchange theorems:
\begin{align*}
 &L\to\infty &&\text{thermodynamic limit},\\
 &a\to0 &&\text{continuum limit},\\
 &t\to\infty &&\text{spectral extraction},\\
 &k\to k_*(a) &&\text{RG iteration to a physical scale},\\
 &A_n\to A &&\text{renormalized observable approximation}.
\end{align*}
The safe route in Theorem~\ref{thm:regcont} takes the regulated limits at each
fixed physical $t$, establishes a continuum positive semigroup estimate for
all $t$, and only then takes $t\to\infty$ inside the spectral lemma.  Directly
choosing $t=t(a)\to\infty$ requires a stronger uniform error theorem.

\section{End-to-end theorem graph}

The proposed proof chain is
\[
\begin{gathered}
 \text{regulated gauge theory}
 \xrightarrow{\text{typed RG}}
 \mathrm{BE}(k_*,\kappa,q)
 \xrightarrow{\text{Thm.~\ref{thm:basin}}}
 \text{UPRR}
 \\
 \xrightarrow{\text{continuum coherence}}
 \text{positive continuum correlations}
 \xrightarrow{\text{OS+density}}
 \spec(H|_{\Omega^\perp})\subset[m,\infty).
\end{gathered}
\]
The middle implications are conditional theorems.  The first arrow to basin
entry and the continuum-coherence arrow are open in the required four-
dimensional strength.

\section{Detailed implementation phases and acceptance gates}

\begin{longtable}{p{.06\textwidth}p{.35\textwidth}p{.20\textwidth}p{.25\textwidth}}
\toprule
Phase & Implemented output & Gate & Current status\\
\midrule
\endhead
0 & exact target, distinctions, C0--C6 hierarchy & no ambiguity in theorem & \Pstatus\\
1 & observable taxonomy and density obligation & state separation proved & \Ostatus\\
2 & representation packages, morphisms, certificates & composition/error theorem & \Pstatus\\
3 & regulator atlas independent of prior Wilson work & one admissible route selected & \Cstatus\\
4 & UPRR definition in physical units & false-positive suite passed & \Pstatus\ as definition\\
5 & positive spectral and regulated transfer theorems & positivity+density supplied & \Cstatus\\
6 & typed gauge-covariant RG architecture & full step estimates & \Ostatus\\
7 & gauge-invariant large-field/topology obligations & summable global control & \Ostatus\\
8 & weighted influence basin & local noncircular criterion & \Pstatus\ abstractly\\
9 & basin-to-rigidity theorem & physical rate and pullback & \Cstatus\\
10 & four-dimensional basin entry & Eq.~\eqref{eq:BE} uniform & \Ostatus\\
11 & continuum Schwinger functions and theory identity & OS, nontrivial, universal & \Ostatus\\
12 & OS reconstruction and spectral extraction & dense core and unique vacuum & \Cstatus\\
13 & restricted-model calibration & massive accepted, false positives rejected & \Rstatus\\
14 & executable finite core & tests pass & \Rstatus\\
15 & frontier audit and claim downgrade & first open implication named & \Pstatus\\
\bottomrule
\end{longtable}

\section{Frontier audit}

\subsection*{What has been proved here}
\begin{itemize}
\item typed representation-error composition;
\item the positive Laplace-support lemma and dense-core spectral transfer;
\item an abstract weighted-influence basin-to-clustering theorem;
\item the conditional logical chain from continuum UPRR and OS density to a
      Hamiltonian gap;
\item counterexamples invalidating finite-volume, lattice-unit, incomplete-
      sector, signed-measure, and stochastic-time shortcuts;
\item finite exact-arithmetic calibration of the spectral and morphism core.
\end{itemize}

\subsection*{What is conditional}
The RG-to-rigidity and continuum-gap theorems are valid under their displayed
hypotheses.  Those hypotheses are not asserted to be known for four-dimensional
pure Yang--Mills.

\subsection*{What remains open}
\begin{enumerate}
\item cutoff- and volume-uniform full-measure entry into a noncircular rigidity
      basin along the asymptotically free trajectory;
\item summable control of large fields, global gauge orbit geometry, and
      topological sectors across all required RG steps;
\item construction, nontriviality, uniqueness, OS regularity, and universality
      of the four-dimensional continuum Schwinger functions;
\item density of the chosen gauge-invariant observable core in the reconstructed
      physical Hilbert space.
\end{enumerate}

Therefore the strongest justified claim is C3, not C6.  The big-picture value
of representational dynamics is that it localizes the unresolved physics: the
problem is not merely to ``find a gap'', but to generate an independently
verifiable infrared contraction and transport it without losing gauge
invariance, positivity, physical scaling, state completeness, or continuum
identity.

\section{Next theorem-level work}

The next productive work is not another rephrasing of the final spectral
statement.  It is to choose one regulator whose conditional block kernels can
be defined globally and attack Eq.~\eqref{eq:BE}.  Three falsifiable branches
are available:
\begin{enumerate}
\item \emph{Influence branch:} prove a weighted conditional influence norm
      below one after a finite physical number of RG steps.
\item \emph{Polymer branch:} derive a convergent full-measure polymer expansion
      at the entry scale whose connected activities yield temporal clustering.
\item \emph{Transfer branch:} prove a gauge-physical transfer-operator
      contraction uniform in volume and cutoff, then compare it to UPRR.
\end{enumerate}
Each branch must include a large-field/topology lemma and a positivity/density
certificate.  Failure of one branch should yield a counterexample or an exact
lower-level obstruction, not be hidden by moving to the next representation.

\begin{thebibliography}{99}
\bibitem{JaffeWitten}
A. Jaffe and E. Witten,
\emph{Quantum Yang--Mills Theory}, Clay Mathematics Institute problem
description,\newline
\url{https://www.claymath.org/wp-content/uploads/2022/06/yangmills.pdf}.

\bibitem{OS1}
K. Osterwalder and R. Schrader,
``Axioms for Euclidean Green's Functions,''
\emph{Communications in Mathematical Physics} \textbf{31} (1973), 83--112,\newline
\url{https://projecteuclid.org/euclid.cmp/1103858969}.

\bibitem{OS2}
K. Osterwalder and R. Schrader,
``Axioms for Euclidean Green's Functions II,''
\emph{Communications in Mathematical Physics} \textbf{42} (1975), 281--305.

\bibitem{BalabanI}
T. Ba\l aban,
``Renormalization Group Approach to Lattice Gauge Field Theories I:
Generation of Effective Actions in a Small Field Approximation and a Coupling
Constant Renormalization in Four Dimensions,''
\emph{Communications in Mathematical Physics} \textbf{109} (1987), 249--301,\newline
\url{https://projecteuclid.org/euclid.cmp/1104116842}.

\bibitem{SZZ}
H. Shen, R. Zhu, and X. Zhu,
``A Stochastic Analysis Approach to Lattice Yang--Mills at Strong Coupling,''
\emph{Communications in Mathematical Physics} \textbf{400} (2023), 805--851,\newline
\url{https://arxiv.org/abs/2204.12737}.

\bibitem{Nissim}
R. Nissim,
``$\mathrm{U}(N)$ Lattice Yang--Mills in the 't Hooft Regime,'' preprint
(2025),\newline
\url{https://arxiv.org/abs/2510.22788}.
\end{thebibliography}

\end{document}
